\section{Open problems}\label{sec:open}

\begin{itemize}
  \item \emph{Determinism.} Replay requires that the same engine, inventory and search order give the same reading
  on every machine. This has not been measured. Where implementations differ, a dispute has no right answer; The
  Graph resolves such cases as a draw \citep{thegraph2026docs}.
  \item \emph{The shape of the reward.} Where the threshold $n_0$ sits and how credit rises above it. A low
  threshold admits small groups; a steep rise favours what is already widely held.
  \item \emph{The value of $p$.} Condition (\ref{eq:ring}) can be met only if the chance that one party posing as
  several Bonfires is caught is known to within a margin, and stays there when few are looking.
  \item \emph{Deciding a judged challenge.} Who votes, what a challenger must post, and what a mistaken slash costs
  an honest party. The argument of Section~\ref{sec:security} assumes mistaken slashes are rare.
  \item \emph{Nodes behind a Bonfire at fault.} Whether, and in what part, a node loses its pledge when the Bonfire
  it serves is found at fault (Section~\ref{sec:bond}).
  \item \emph{Parameters.} The ceiling rate $c$, the pledge $q$, the minimum bond $m$, the challenger's fraction
  $\varphi$, the lineage share $\lambda$, the length of the lock, the rate of spot checks and what they are paid,
  and whether the protocol bounds the share $s$.
  \item \emph{Necessity and granularity.} A reading may survive the removal of either of two constructions and fail
  without both; the test then credits neither. Attribution rules of the Shapley kind would handle this but are
  vulnerable when a participant can appear under more than one name \citep{ohta2008anonymity}.
  \item \emph{Comparing meanings.} The copy rule asks whether an earlier construction gives the same reading. A copy
  that renames the symbols of the meaning gives a different one, and a test of equivalence is wanted.
  \item \emph{Registration in bulk.} Priority belongs to the first to register, and constructions can be generated
  and registered in any number ahead of the communities that would have made them. What registration should cost
  is open.
  \item \emph{Declaring parents.} A construction that declares a parent gives up a share of its credit. A test of
  derivation that anyone can run is wanted.
  \item \emph{The base grammar.} It earns nothing, and whoever changes it can make any construction unnecessary.
  Who maintains it, and how that work is paid for, is open.
  \item \emph{What a node costs.} A node is paid a share of issuance. Whether that covers the computation it
  supplies, and who pays where it does not, is open.
  \item \emph{Disclosure.} A challenge shows the turns of one encounter to the challenger.
  \item \emph{Payment from use.} The protocol is funded by issuance alone. Whether payment from those who use the
  work of agents further down the supply chain should stand behind it is open.
  \item \emph{Learning without teaching.} Constructions enter a grammar today by being taught. Deriving them from a
  group's discourse is a proposed experiment of the companion paper \citep{vera2026constructions}.
\end{itemize}

Most of these would be settled by simulation before deployment: a naming game among Bonfires that hold different
grammars, with lookups, adoptions and lineage recorded, and the same game with a party that poses as several
Bonfires, a group that agrees to hold one another's constructions, and a party that borrows stake.
